% !PRINTED 387
\weberpage{387}
%<<CONT
\weq{7}{
\begin{aligned}
&\left(\frac{a}{l}\right),\ \left(\frac{a}{l'}\right), \dots && m \equiv 3 \pmod 4, \\
&\left(\frac{-1}{a}\right),\ \left(\frac{a}{l}\right),\ \left(\frac{a}{l'}\right), \dots && m \equiv 1 \pmod 4, \\
& && m \equiv 4 \pmod 8, \\
&\left(\frac{2}{a}\right),\ \left(\frac{a}{l}\right),\ \left(\frac{a}{l'}\right), \dots && m \equiv 6 \pmod 8, \\
&\left(\frac{-2}{a}\right),\ \left(\frac{a}{l}\right),\ \left(\frac{a}{l'}\right), \dots && m \equiv 2 \pmod 8, \\
&\left(\frac{-1}{a}\right),\ \left(\frac{2}{a}\right),\ \left(\frac{a}{l}\right),\ \left(\frac{a}{l'}\right), \dots && m \equiv 0 \pmod 8.
\end{aligned}
}
The number of these fundamental characters is, as one sees, equal to
the number $\lambda$ determined more closely in \S\,104, 3., and the number of
sign combinations is equal to $2^\lambda$.

The dependence among these characters may be represented as
follows.

If
\[
m = n^2 m', \quad D = Q^2\Delta = -4n^2m'
\]
and $n^2$ is the greatest square number dividing $m$, then
\weq{8}{
\begin{aligned}
&\Delta \equiv -m', && \text{if } m' \equiv 3 \pmod 4, \\
&\Delta \equiv -4m' && \text{,,}\quad m' \equiv 1, 2 \pmod 4;
\end{aligned}
}
in the first case $m'$ is odd and not equal to 1, hence equal to
one of the primes $l$ or equal to a product of several
of them, and from
\[
\left(\frac{\Delta}{a}\right) = \left(\frac{-m'}{a}\right) = \left(\frac{a}{m'}\right) = 1,
\]
in which the law of reciprocity of quadratic residues has been
applied, a relation between the characters results.

In the second case, for odd $m'$,
\[
\left(\frac{-m'}{a}\right) = \left(\frac{-1}{a}\right)\left(\frac{a}{m'}\right) = 1,
\]
and this is again a product of several of the characters (7),
and if finally $m' = 2m''$ is even, then
\[
\left(\frac{-m'}{a}\right) = \left(\frac{\pm 2}{a}\right)\left(\frac{a}{m''}\right),
\]
where the sign $\pm$ is so determined that $\pm m'' \equiv 1 \pmod 4$.
Thus again the number of genera is at most $= 2^{\lambda-1}$.
